% =============================
% ภาคผนวก
% =============================

% =============================
% ภาคผนวก ก
% =============================
\setcounter{chapter}{-1}
\appendiceschapter{คู่มือการติดตั้งและการใช้งานระบบ}

\section*{ข้อกำหนดเบื้องต้นของสภาพแวดล้อม (Prerequisites)}

\begin{itemize}
    \item ระบบปฏิบัติการ: Windows 10/11, macOS หรือ Linux (Ubuntu 20.04+)
    \item Python เวอร์ชัน 3.11 หรือใหม่กว่า
    \item เครื่องมือจัดการสภาพแวดล้อม \texttt{uv} (Fast Python Package Installer)
    \item การเชื่อมต่ออินเทอร์เน็ตสำหรับดาวน์โหลด Dependencies และเรียกใช้ KKU Gen.AI API
\end{itemize}

\section*{ขั้นตอนการติดตั้งระบบ (Installation Steps)}

\begin{enumerate}[label=\arabic*.]
    \item โคลน Repository โครงการจาก GitHub:
    \begin{lstlisting}[language=bash]
git clone https://github.com/wittgenstein-byte/End-to-End_News_Collector-Summarizer.git
cd End-to-End_News_Collector-Summarizer
    \end{lstlisting}

    \item ติดตั้ง Dependencies ทั้งหมดผ่าน \texttt{uv}:
    \begin{lstlisting}[language=bash]
uv sync
    \end{lstlisting}

    \item ติดตั้ง Playwright Chromium Headless Browser:
    \begin{lstlisting}[language=bash]
uv run playwright install chromium
    \end{lstlisting}

    \item กำหนดค่า Environment Variables โดยคัดลอกไฟล์ตัวอย่าง:
    \begin{lstlisting}[language=bash]
copy backend\.env.example backend\.env
    \end{lstlisting}
    จากนั้นแก้ไขไฟล์ \texttt{backend/.env} โดยกรอก API Key สำหรับ LLM Service:
    \begin{lstlisting}[language=bash]
LLM_API=your_kku_genai_api_key_here
    \end{lstlisting}

    \item เริ่มต้นทำงานเซิร์ฟเวอร์ Backend:
    \begin{lstlisting}[language=bash]
uv run python backend/main.py
    \end{lstlisting}

    \item เปิดเว็บเบราว์เซอร์และเข้าสู่ระบบที่: \url{http://localhost:5000}
\end{enumerate}

% =============================
% ภาคผนวก ข
% =============================
\appendiceschapter{ข้อกำหนด REST API และ WebSocket Events}

\section*{รายการ REST API Endpoints}

\begin{table}[H]
\centering
\caption{รายละเอียด REST API Endpoints ของระบบ}
\label{tab:api_specs}
\resizebox{\textwidth}{!}{%
\begin{tabular}{|l|l|p{5cm}|p{4.5cm}|}
\hline
\textbf{Method} & \textbf{Endpoint} & \textbf{คำอธิบาย} & \textbf{Parameters / Response} \\
\hline
\texttt{GET} & \texttt{/api/news} & ดึงรายการข่าวแบบแบ่งหน้า & \texttt{page}, \texttt{category}, \texttt{search} \\ \hline
\texttt{GET} & \texttt{/api/news/\{id\}} & ดึงรายละเอียดข่าวรายชิ้นพร้อมบทสรุป & คืนค่า JSON \texttt{NewsArticle} \\ \hline
\texttt{GET} & \texttt{/api/categories} & ดึงรายการหมวดหมู่ข่าวทั้งหมด & คืนค่า List ของ Category Metadata \\ \hline
\texttt{GET} & \texttt{/api/trending} & ดึงรายการข่าวที่เป็นกระแสความนิยม & คืนค่า List ของ \texttt{TrendingArticle} \\ \hline
\texttt{POST} & \texttt{/api/engagement} & บันทึกสถิติการปฏิสัมพันธ์ของผู้ใช้ & รับ \texttt{article\_id}, \texttt{action\_type} \\ \hline
\texttt{POST} & \texttt{/api/collect} & สั่งให้ระบบเริ่มรอบการสเกรปข่าวทันที & คืนค่าสถานะ \texttt{status: started} \\ \hline
\texttt{GET} & \texttt{/health} & ตรวจสอบสถานะความพร้อมของเซิร์ฟเวอร์ & คืนค่าสถานะ Storage และ Playwright \\
\hline
\end{tabular}%
}
\end{table}

\section*{รายการ WebSocket Events (Socket.IO)}

\begin{table}[H]
\centering
\caption{รายละเอียด WebSocket Events สำหรับการสื่อสารแบบ Real-time}
\label{tab:ws_events}
\resizebox{\textwidth}{!}{%
\begin{tabular}{|l|l|p{8cm}|}
\hline
\textbf{Event Name} & \textbf{ทิศทาง} & \textbf{คำอธิบายและข้อมูลที่ส่ง (Payload)} \\
\hline
\texttt{news:new} & Server $\rightarrow$ Client & ส่งข้อมูลข่าวใหม่ที่เพิ่งผ่านการประมวลผลเสร็จสิ้น เพื่อแสดงผลบนหน้าจอทันที \\ \hline
\texttt{scraper:status} & Server $\rightarrow$ Client & แจ้งสถานะความคืบหน้าของกระบวนการสเกรปข่าว (เช่น เริ่มต้น, จำนวนที่ดึงได้, เสร็จสิ้น) \\ \hline
\texttt{news:updated} & Server $\rightarrow$ Client & ส่งข้อมูลการอัปเดตบทความ เช่น เมื่อการสรุป LLM เสร็จสมบูรณ์ \\ \hline
\texttt{connect} / \texttt{disconnect} & ทั้งสองทิศทาง & ตรวจจับการเชื่อมต่อและการตัดการเชื่อมต่อของผู้ใช้ \\
\hline
\end{tabular}%
}
\end{table}
